\documentclass[a4paper,twoside]{article}

\usepackage{morewrites}

\usepackage[svgnames,dvipsnames,x11names]{xcolor}
\usepackage{graphicx}
\usepackage[english]{babel}
\usepackage[colorlinks=true, linkcolor=NavyBlue, urlcolor=NavyBlue, citecolor=NavyBlue]{hyperref}
\usepackage[a4paper]{geometry}
\usepackage{ts-fonts}
\usepackage{amsmath}
\usepackage{float}
\usepackage{seqsplit}
\usepackage{ts-callouts}
\usepackage{ts-code}

\usepackage{lastpage}
\usepackage{refcount}
\makeatletter
\def\ps@plain{
  \let\@oddhead\@empty
  \let\@evenhead\@empty
  \def\@oddfoot{
    \hfil
    \ifnum\getpagerefnumber{LastPage}>1\relax
      \thepage
    \fi
    \hfil}
  \let\@evenfoot\@oddfoot
}
\makeatother

\title{Output backends}
\author{}\date{}

\begin{document}

\maketitle

\thispagestyle{plain}
TeXSmith lowers every document into a typed intermediate representation (IR)
and then emits a backend from that IR:

\begin{code}{text}{}{}
\begin{Verbatim}[commandchars=\\\{\},breaklines, breakanywhere, commandchars=\\\{\}]
read(HTML) → IR (texsmith.ir) → write(IR) → LaTeX | Typst
\end{Verbatim}
\end{code}
Because both backends consume the \textbf{same} IR (the readers and the IR are
untouched), you choose the output format with a single flag:

\begin{code}{bash}{}{}
\begin{Verbatim}[commandchars=\\\{\},breaklines, breakanywhere, commandchars=\\\{\}]
texsmith\PY{+w}{ }doc.md\PY{+w}{ }\PYZhy{}tarticle\PY{+w}{ }\PYZhy{}\PYZhy{}format\PY{+w}{ }latex\PY{+w}{   }\PY{c+c1}{\PYZsh{} default}
texsmith\PY{+w}{ }doc.md\PY{+w}{ }\PYZhy{}tarticle\PY{+w}{ }\PYZhy{}\PYZhy{}format\PY{+w}{ }typst
\end{Verbatim}
\end{code}
\texttt{-\allowbreak{}-\allowbreak{}format} is case-insensitive and accepts \texttt{latex} (default) or \texttt{typst}.

\section{LaTeX backend (default)}\label{latex-backend-default}

The \LaTeX{} backend (\texttt{texsmith.writers.latex}) is the full-featured path: it
drives the template/fragment runtime, fonts and script matching, glossary and
index engines, asset transformers, and compiles through a TeX engine
(Tectonic by default, or \texttt{latexmk} with \texttt{-\allowbreak{}-\allowbreak{}engine lualatex} / \texttt{-\allowbreak{}-\allowbreak{}engine
xelatex}). This is the backend assumed throughout most of this documentation.

\section{Typst backend}\label{typst-backend}

\begin{callout}[callout warning]{Experimental}
The Typst backend is \textbf{experimental}. The \LaTeX{} backend remains the
default and the supported path; Typst covers a subset of it, leans on the
third-party \texttt{mitex} package for math (which can fail on some constructs ---
see below), and approximates or drops a few constructs. Use it for new,
Typst-first documents and check the output; do not assume \LaTeX{}-level
fidelity yet.
\end{callout}
\href{https://typst.app}{Typst} is a modern, Rust-based typesetting system. The
Typst backend (\texttt{texsmith.writers.typst}) emits a compilable \texttt{.typ} source from
the same IR. It is a lean, Typst-native path and intentionally covers a
\textbf{subset} of the \LaTeX{} backend (see \hyperref[scope-and-limitations]{Scope}).

\begin{code}{bash}{}{}
\begin{Verbatim}[commandchars=\\\{\},breaklines, breakanywhere, commandchars=\\\{\}]
\PY{c+c1}{\PYZsh{} Emit a .typ file (no compiler required)}
texsmith\PY{+w}{ }doc.md\PY{+w}{ }\PYZhy{}tarticle\PY{+w}{ }\PYZhy{}\PYZhy{}format\PY{+w}{ }typst\PY{+w}{ }\PYZhy{}\PYZhy{}output\PY{+w}{ }build/

\PY{c+c1}{\PYZsh{} Emit and compile to PDF}
texsmith\PY{+w}{ }doc.md\PY{+w}{ }\PYZhy{}tarticle\PY{+w}{ }\PYZhy{}\PYZhy{}format\PY{+w}{ }typst\PY{+w}{ }\PYZhy{}\PYZhy{}build\PY{+w}{ }\PYZhy{}\PYZhy{}output\PY{+w}{ }build/
\end{Verbatim}
\end{code}
Without a template, the body is wrapped in a minimal standalone preamble. With
the \textbf{article} or \textbf{book} template, the body is wrapped in the template's
Typst scaffolding (title, authors, date, abstract, table of contents,
sectioning, and a native Typst bibliography).

\begin{callout}[callout note]{Typst output and \LaTeX{}-only flags}
\texttt{-\allowbreak{}-\allowbreak{}format typst} cannot be combined with \texttt{-\allowbreak{}-\allowbreak{}html}, \texttt{-\allowbreak{}-\allowbreak{}template-\allowbreak{}info}, or
\texttt{-\allowbreak{}-\allowbreak{}template-\allowbreak{}scaffold}: those inspect the \LaTeX{} fragment/engine machinery,
which the Typst path does not use.
\end{callout}
\subsection{Installing a Typst compiler}\label{installing-a-typst-compiler}

Emitting the \texttt{.typ} source never needs a compiler; compilation does. Two paths
are supported, tried in this order:

\textbf{Embedded compiler (pure pip)}\par

The \texttt{typst} PyPI package embeds the Rust compiler, so nothing needs to be
on your \texttt{PATH}:

\begin{code}{bash}{}{}
\begin{Verbatim}[commandchars=\\\{\},breaklines, breakanywhere, commandchars=\\\{\}]
pip\PY{+w}{ }install\PY{+w}{ }\PY{l+s+s2}{\PYZdq{}texsmith[typst]\PYZdq{}}
\PY{c+c1}{\PYZsh{} or}
uv\PY{+w}{ }pip\PY{+w}{ }install\PY{+w}{ }\PYZhy{}e\PY{+w}{ }\PY{l+s+s2}{\PYZdq{}.[typst]\PYZdq{}}
\PY{c+c1}{\PYZsh{} or}
uv\PY{+w}{ }tool\PY{+w}{ }install\PY{+w}{ }\PY{l+s+s2}{\PYZdq{}texsmith[typst]\PYZdq{}}
\end{Verbatim}
\end{code}
\textbf{System binary}\par

Install \texttt{typst} on your \texttt{PATH} and TeXSmith detects it automatically as a
fallback:

\begin{code}{bash}{}{}
\begin{Verbatim}[commandchars=\\\{\},breaklines, breakanywhere, commandchars=\\\{\}]
brew\PY{+w}{ }install\PY{+w}{ }typst\PY{+w}{          }\PY{c+c1}{\PYZsh{} Homebrew}
cargo\PY{+w}{ }install\PY{+w}{ }typst\PYZhy{}cli\PY{+w}{      }\PY{c+c1}{\PYZsh{} Cargo}
\PY{c+c1}{\PYZsh{} or download a release from https://github.com/typst/typst/releases}
\end{Verbatim}
\end{code}
\begin{callout}[callout warning]{Prefer the system binary for recent math}
The compiler embedded in the \texttt{typst} PyPI package can lag behind recent
Typst releases. Math is rendered through the \texttt{mitex} package (a \LaTeX{}-math
renderer for Typst); documents that rely on a recent \texttt{mitex} may \textbf{fail}
with the embedded compiler where an up-to-date system binary only emits a
warning. For math-heavy documents, prefer the system binary.
\end{callout}
If no compiler is available, the \texttt{.typ} is still written and TeXSmith reports
that compilation was skipped, with an actionable install hint.

\subsection{Scope and limitations}\label{scope-and-limitations}

The Typst backend covers a real templated document (article and book) and the
common Markdown/HTML constructs:

\begin{itemize}
\item{} Document structure, paragraphs, and headings.
\item{} Inline emphasis (emphasis, strong, strikeout, underline, highlight, small
  caps, sub/superscript, quotes), inline code, and links.
\item{} Code blocks, block quotes, bullet and ordered lists, definition lists,
  admonitions, horizontal rules.
\item{} Images, figures, simple (GFM) tables, and rich (\texttt{yaml} / \texttt{data-\allowbreak{}ts}) tables
  rebuilt as a native Typst \texttt{\#table}.
\item{} Footnotes, TeX logos, keystrokes, progress bars.
\item{} Math (rendered through the Typst \texttt{mitex} package) and equation
  cross-references (\texttt{\#ref}).
\item{} Native citations (\texttt{\#cite}) resolved against the bibliography collection the
  CLI builds (\texttt{.bib} files plus inline DOI), wired through \texttt{\#bibliography(...)}.

\end{itemize}
Approximated or dropped on the Typst path (no exact equivalent --- the content is
preserved as closely as possible rather than producing wrong output):

\begin{itemize}
\item{} \texttt{MarginNote} renders as a \texttt{\#footnote} (the \texttt{side} hint has no counterpart).
\item{} Index entries (\texttt{<abbr>} markers) keep their visible text but emit no index
  marker --- Typst has no \texttt{makeindex} equivalent here.
\item{} Remote (\texttt{http(s)://}, \texttt{data:}) or unresolved local images are dropped to keep
  the document compilable.
\item{} The LaTeX-only glossary/index engines and the fragment system are not used on
  the Typst path.

\end{itemize}
Any IR node that still has \textbf{no} Typst emitter raises an explicit, localised
\texttt{TypstWriteError} (naming the node and the backend) rather than emitting wrong
output. For anything beyond this subset, use the default \LaTeX{} backend.

\end{document}
